% ============================================================
% The Collision Transform in the Critical Strip
% ============================================================

\documentclass[12pt]{amsart}

\usepackage{amsmath,amssymb,amsthm}
\usepackage{booktabs}
\usepackage{microtype}
\usepackage{url}

\newtheorem{theorem}{Theorem}
\newtheorem{lemma}[theorem]{Lemma}
\newtheorem{proposition}[theorem]{Proposition}
\newtheorem{corollary}[theorem]{Corollary}
\newtheorem{conjecture}[theorem]{Conjecture}
\theoremstyle{definition}
\newtheorem{definition}[theorem]{Definition}
\theoremstyle{remark}
\newtheorem{remark}[theorem]{Remark}

\title{The Collision Transform\\and the Critical Strip}

\author{Alexander S.\ Petty}

\email{alexander.petty@gmail.com}
\date{June 2025}
\subjclass[2020]{11A63, 11M06, 11M26}

\begin{document}
\begin{abstract}
The centered collision sum
$F^{\circ}(s) = \sum_p S^{\circ}_p / p^s$
converges at $s = 1$ (proved in the preceding paper).
This paper investigates whether convergence extends
below $s = 1$ into the critical strip.

The collision transform writes $F^{\circ}(s)$ as a
finite linear combination of character sums
$P(s, \chi) = \sum_p \chi(p)/p^s$, linked to
$\log L(s, \chi)$ for $L$-functions modulo
$b^{\ell+1}$. Antisymmetry restricts the sum to odd
characters. Conditional on a zero-free hypothesis for
the relevant $L$-functions, $F^{\circ}(s)$ converges
for real $s > \sigma_0$, where $\sigma_0$ depends on
the zero-free region.

Computation suggests convergence persists to
$s \approx 0.5$ in base~$3$ and to $s \approx 0.6$
in bases $5$, $7$, $10$, $12$. These observations
are consistent with, but do not prove, convergence
in the strip.
\end{abstract}
\maketitle

% ============================================================
\section{Introduction}

This paper investigates the collision transform below $s = 1$.

The preceding paper~\cite{paper16} proved that the centered
collision fluctuation sum
\[
F^{\circ}(1, \ell)
= \sum_{\chi \ne \chi_0}
\hat{S}_{\ell}(\chi) \sum_p \frac{\chi(p)}{p}
\]
converges at $s = 1$. The proof used the finite determination
theorem (the collision deviation $S_{\ell}(p)$ depends only on
$p \bmod b^{\ell+1}$), the collision transform (decomposition
over Dirichlet characters modulo $b^{\ell+1}$), and Mertens'
theorem for arithmetic progressions.

The present paper investigates whether convergence extends
below $s = 1$.

% ============================================================
\section{The Character Decomposition Below $s = 1$}

For $\operatorname{Re}(s) > 1$, the prime character
sum~\cite{davenport,iwaniec-kowalski}
\[
P(s, \chi) = \sum_p \frac{\chi(p)}{p^s}
\]
is related to the Dirichlet $L$-function by
\begin{equation}\label{eq:P-and-L}
P(s, \chi) = \log L(s, \chi) - H(s, \chi),
\end{equation}
where the higher-power sum
$H(s,\chi) = \sum_p \sum_{k\ge2} \chi(p)^k/(k\,p^{ks})$
converges absolutely for $\operatorname{Re}(s)>1/2$.

The centered fluctuation sum is therefore
\[
F^{\circ}(s) = -\sum_{\chi \ne \chi_0}
\hat{S}(\chi)\, \log L(s, \chi)
+ \sum_{\chi \ne \chi_0}
\hat{S}(\chi)\, H(s, \chi).
\]
The second sum converges for $\operatorname{Re}(s) > 1/2$.
The convergence of $F^{\circ}(s)$ at $s = \sigma$ is
therefore controlled by the first sum: the weighted
combination of $\log L(s, \chi)$ values at
$s = \sigma$.

If $L(\sigma, \chi)$ has a zero at some $\rho$ with
$\operatorname{Re}(\rho) > \sigma$, then
$\log L(s, \chi)$ has a logarithmic singularity at
$\rho$. Convergence of $F^{\circ}(\sigma)$ despite
such singularities would require cancellation across
characters: the contributions from zeros of different
$L$-functions, weighted by $\hat{S}(\chi)$, must cancel.

% ============================================================
\section{Computational Evidence}

\begin{table}[h]
\centering
\begin{tabular}{rrrrrrr}
\toprule
$s$ & base $3$ & base $5$ & base $7$ & base $10$ & base $12$ \\
\midrule
$1.0$ & $-0.151$ & $-0.187$ & $-0.137$ & $-0.199$ & $-0.278$ \\
$0.9$ & $-0.168$ & $-0.230$ & $-0.186$ & $-0.282$ & $-0.417$ \\
$0.8$ & $-0.186$ & $-0.277$ & $-0.250$ & $-0.402$ & $-0.624$ \\
$0.7$ & $-0.205$ & $-0.322$ & $-0.337$ & $-0.577$ & $-0.926$ \\
$0.6$ & $-0.227$ & $-0.347$ & $-0.457$ & $-0.835$ & $-1.346$ \\
$0.5$ & $-0.257$ & $-0.312$ & $-0.636$ & $-1.219$ & $-1.880$ \\
\bottomrule
\end{tabular}
\caption{$F^{\circ}(s)$ at lag $\ell = 1$ across bases,
using $18{,}000$--$78{,}000$ primes per base.}
\label{tab:values}
\end{table}

\begin{table}[h]
\centering
\begin{tabular}{rrrrrrr}
\toprule
$s$ & base $3$ & base $5$ & base $7$ & base $10$ & base $12$ \\
\midrule
$1.0$ & $0.000$ & $0.000$ & $0.000$ & $0.000$ & $0.000$ \\
$0.9$ & $0.000$ & $0.001$ & $0.001$ & $0.000$ & $0.001$ \\
$0.8$ & $0.000$ & $0.004$ & $0.002$ & $0.001$ & $0.005$ \\
$0.7$ & $0.000$ & $0.013$ & $0.007$ & $0.004$ & $0.017$ \\
$0.6$ & $0.001$ & $0.042$ & $0.027$ & $0.016$ & $0.058$ \\
$0.5$ & $0.003$ & $0.135$ & $0.099$ & $0.065$ & $0.205$ \\
\bottomrule
\end{tabular}
\caption{Drift: absolute change in $F^{\circ}(s)$ between
the first and second halves of the prime range. Small drift
indicates convergence. Drift less than $0.01$ is consistent
with stability.}
\label{tab:drift}
\end{table}

The data reveals a base-dependent penetration pattern:

\begin{itemize}
\item \textbf{Base~$3$}: drift $< 0.003$ at all $s$ down to
$s = 0.5$. The centered sum appears to converge on the
critical line.

\item \textbf{Base~$5$}: stable to $s = 0.7$. Drift grows at
$s = 0.6$ and $s = 0.5$.

\item \textbf{Bases~$7$, $10$}: stable to $s = 0.7$. Moderate
drift at $s = 0.6$. Growing drift at $s = 0.5$.

\item \textbf{Base~$12$}: stable to $s = 0.7$. Significant
drift at $s = 0.6$ and $s = 0.5$.
\end{itemize}

% ============================================================
\section{The Penetration Depth}

\begin{definition}
The \emph{penetration depth} of the collision transform at
base $b$ and lag $\ell$ is
\[
\sigma_0(b, \ell) = \inf\{\sigma > 0 :
F^{\circ}(\sigma, \ell) \text{ converges}\}.
\]
\end{definition}

The data suggests $\sigma_0(3, 1) \le 0.5$ and
$\sigma_0(b, 1) \le 0.7$ for all bases tested.

The base dependence of the apparent penetration depth
may reflect the number of primes required rather than
a genuine difference in $\sigma_0$. Base~$3$ has
$\phi(9) = 6$ characters modulo $b^2 = 9$, while
base~$12$ has $\phi(144) = 48$. Cancellation among
more characters requires more primes to manifest.

\begin{conjecture}[Universal penetration]\label{conj:penetration}
For every base $b \ge 2$ and lag $\ell \ge 1$:
$\sigma_0(b, \ell) = 1/2$.
\end{conjecture}

This conjecture is consistent with GRH. By
equation~\eqref{eq:P-and-L}, the convergence of
$F^{\circ}(s)$ for $\operatorname{Re}(s) > 1/2$ is
equivalent to the convergence of
$\sum_{\chi \ne \chi_0} \hat{S}(\chi)\, \log L(s, \chi)$
for $\operatorname{Re}(s) > 1/2$. Under GRH, the
individual $P(s, \chi)$ converge for
$\operatorname{Re}(s) > 1/2$, and the finite sum
$F^{\circ}(s)$ inherits this convergence.

% ============================================================
\section{The Antisymmetry Theorem}

\begin{theorem}[Odd-character restriction]\label{thm:odd}
The centered collision transform satisfies
$\hat{S}^{\circ}(\chi) = 0$ for every even Dirichlet
character $\chi$ modulo $b^{\ell+1}$.
\end{theorem}

\begin{proof}
By the reflection identity~\cite{paper16},
$S(a) + S(m{-}a) = -1$ for $m = b^{\ell+1}$. The map
$a \mapsto m - a$ sends spectral class $R$ to class
$b - 2 - R$, with class means satisfying
$\overline{S}_R + \overline{S}_{b-2-R} = -1$. Therefore
$S^{\circ}(a) + S^{\circ}(m{-}a) = 0$.

For any character $\chi$, $\chi(m{-}a) = \chi(-1)\chi(a)$.
If $\chi$ is even ($\chi(-1) = 1$):
\[
\hat{S}^{\circ}(\chi) = \frac{1}{\phi(m)}
\sum_a S^{\circ}(a)\, \overline{\chi}(a)
= \frac{1}{\phi(m)}
\sum_a S^{\circ}(a)\, \overline{\chi}(m{-}a)
= -\hat{S}^{\circ}(\chi),
\]
using $S^{\circ}(m{-}a) = -S^{\circ}(a)$ and
$\overline{\chi}(m{-}a) = \overline{\chi}(a)$.
Therefore $\hat{S}^{\circ}(\chi) = 0$.
\end{proof}

\begin{corollary}
The centered fluctuation sum decomposes as
\[
F^{\circ}(s) = \sum_{\substack{\chi \bmod m \\
\chi(-1) = -1}} \hat{S}^{\circ}(\chi)\, P(s, \chi).
\]
Only odd $L$-functions control the penetration depth.
\end{corollary}

% ============================================================
\section{Positivity of Character Products}

At the small bases displayed below, each odd character
$\chi$ contributing to $F^{\circ}$ has
\[
\operatorname{Re}\bigl(\hat{S}^{\circ}(\chi)\, P(1/2, \chi)\bigr)\ge 0.
\]
This is only a finite observation, not a general positivity
law.

\begin{table}[h]
\centering
\begin{tabular}{rrrrr}
\toprule
base & odd $\chi$ & $|\hat{S}^{\circ}|$ & $|P(1/2)|$
& $\operatorname{Re}(\hat{S}^{\circ} \cdot P)$ \\
\midrule
$3$ & $j = 1$ & $0.667$ & $0.469$ & $+0.309$ \\
$3$ & $j = 5$ & $0.667$ & $0.469$ & $+0.309$ \\
$5$ & $j = 3$ & $0.828$ & $0.520$ & $+0.368$ \\
$5$ & $j = 17$ & $0.828$ & $0.520$ & $+0.368$ \\
\bottomrule
\end{tabular}
\caption{Character products at selected small bases.
All displayed real parts are non-negative. Later
computation at base~$10$ shows that mixed signs occur
in general.}
\label{tab:positivity}
\end{table}

Later computation shows positivity fails in general:
at base~$10$, $7$ of $20$ products have negative real
part. Convergence of $F^{\circ}(s)$ below $s = 1$
therefore depends on cancellation between terms, not
on individual term convergence.

\begin{remark}[Conditional penetration]
The following conditional statement illustrates one
possible route from the collision transform to GRH,
though the required positivity hypothesis has since
been shown to fail in general.

If the products
$\hat{S}^{\circ}(\chi) \cdot P(s, \chi)$ had
non-negative real part for all odd $\chi$ and all
$s \ge 1/2$, then:
\[
F^{\circ}(1/2) \text{ converges}
\;\Longleftrightarrow\;
L(s, \chi) \ne 0 \;\text{for}\;
\operatorname{Re}(s) > \tfrac{1}{2},\;
\chi \in \mathcal{X}^{-}.
\]
Since positivity fails, the reverse direction of this
argument does not apply. Convergence of $F^{\circ}(s)$
below $s = 1$ may occur through cancellation between
terms of mixed sign, without requiring each individual
$L$-function to be zero-free. The relationship between
$F^{\circ}$ and GRH remains open.
\end{remark}

% ============================================================
\section{The Aperiodic Residual}

The finite determination theorem shows that
$S_{\ell}(p)$ depends only on $p \bmod b^{\ell+1}$. The
centered version $S^{\circ}$ removes the class mean. A
deeper question: does the residual $S^{\circ}$ carry
additional modular structure at higher resolution?

In the tested lag-$1$ cases, computation shows it does
not. For each lag-$1$ class
$a \bmod b^2$, the sub-classes $a' \bmod b^3$ satisfying
$a' \equiv a \pmod{b^2}$ all have identical mean $S$
values. For these examples, the lag-$1$ residual carries
no additional modular information at lag~$2$.

\begin{table}[h]
\centering
\begin{tabular}{rrrrr}
\toprule
base & lag-$1$ mod & lag-$2$ mod & sub-classes & residual \\
\midrule
$3$ & $9$ & $27$ & $3$ per class & $0$ everywhere \\
$5$ & $25$ & $125$ & $5$ per class & $0$ everywhere \\
\bottomrule
\end{tabular}
\caption{Tested lag-$1$ towers in bases $3$ and $5$.
In these examples, the lag-$1$ residual does not split
further at lag~$2$.}
\label{tab:tower}
\end{table}

These examples suggest that the centered deviation
$S^{\circ}$ behaves aperiodically after the lag-$1$ class
mean is removed: in the tested towers, it is not refined
by passing from $b^2$ to $b^3$. The stamp at each lag may
capture all of the modular structure visible at that lag,
with the remainder behaving as a non-periodic fluctuation.

The convergence of $F^{\circ}(s)$ at $s = \sigma$ is
therefore, at least heuristically, a statement about
whether this residual cancels at weighting $1/p^{\sigma}$.
At $s = 1$, the centered sum cancels by the collision
transform (proved). Whether convergence extends below
$s = 1$ depends on cancellation between terms of mixed
sign and remains open.

% ============================================================
\section{The Base Sweep}

The preceding sections analyze the collision transform at a
single base. A different computational perspective emerges
from sweeping across a finite family of bases. For a fixed
prime~$p$ and a finite tested set $\mathcal{B}$, define the
base-swept resistance
\[
\mathcal{R}_{\mathcal{B},w}(p) = \sum_{b \in \mathcal{B}} w(b)\, S_1(p, b),
\]
where $w(b)$ is a chosen weight on $\mathcal{B}$. This
aggregates the collision deviation of a single prime across
the tested bases simultaneously.

In the tested sweeps, $\mathcal{R}_{\mathcal{B},w}(p)$
splits by $p \bmod 3$: primes $\equiv 1 \pmod{3}$ carry
approximately $2/3$ of the total resistance and primes
$\equiv 2 \pmod{3}$ carry approximately $1/3$. This ratio
is stable across the tested weights $w = 1/b$, $w = 1$,
$w = 1/\sqrt{b}$, and $w = 1/b^2$, but no general
invariance theorem is claimed here.

The mod-$3$ structure is not visible from any single base.
In these experiments, it emerges only from the collective
instrument.

% ============================================================
\section{The Mod-$3$ Polarity}

In the tested cases, the centered residual $S^{\circ}$
shows no modular structure in the base direction: the
lag tower appears flat (Section~\ref{tab:tower}). But
it has structure in the \emph{prime} direction.
Its mean over primes $p \equiv 1 \pmod{3}$ differs from
its mean over primes $p \equiv 2 \pmod{3}$. In base~$10$
at lag~$1$, the centered values $S^{\circ}(a)$ for units
$a \equiv 1 \pmod{3}$ average $-1.30$, while those for
$a \equiv 2 \pmod{3}$ average~$0$.

In the tested base sweeps, this asymmetry is stable
under the reweightings described in the previous
section.  A base-swept resistance
\[
\textstyle\sum_{b \in \mathcal{B}} w(b)\, S_1(p, b)
\]
splits by $p \bmod 3$: primes $\equiv 1 \pmod{3}$
carry approximately $2/3$ of the total resistance
and primes $\equiv 2 \pmod{3}$ carry approximately
$1/3$.

Despite this asymmetry, the antisymmetry
$S^{\circ}(a) = -S^{\circ}(m{-}a)$ is compatible with the
mod-$3$ structure because the reflection $a \mapsto m - a$
respects the mod-$3$ pairing (in base~$10$:
$m = 100 \equiv 1 \pmod{3}$).

However, in the computations reported here, removing the
mod-$3$ bias (doubly centering by both spectral class and
$p \bmod 3$) destroys the apparent penetration depth. The
doubly centered sum $F^{\circ\circ}$ still exhibits the
same antisymmetry, but in the tested data it grows rapidly
for $s < 1$.

\begin{table}[h]
\centering
\begin{tabular}{rrr}
\toprule
$s$ & $F^{\circ}$ & $F^{\circ\circ}$ \\
\midrule
$1.0$ & $0.08$ & $0.77$ \\
$0.7$ & $0.27$ & $12.3$ \\
$0.5$ & $0.46$ & $111$ \\
\bottomrule
\end{tabular}
\caption{Removing the mod-$3$ component in the tested
base-$3$ data leaves the $s=1$ value bounded but produces
rapid growth below $s = 1$. This suggests that the
mod-$3$ polarity is relevant to penetration depth.}
\end{table}

The data suggests that the mod-$3$ structure is not needed
for convergence at $s = 1$, but may matter for any deeper
penetration below $s = 1$. At present this is an observed
feature of the computations, not a theorem.

The fraction $2/3$ that governs this distribution is
the alignment limit of the prime~$3$ from the first
paper in this series~\cite{paper1}: the constant
$(p{-}1)/p$ evaluated at the prime the golden ratio
selects. Its reappearance here, as the observed fraction
of base-swept resistance carried by primes
$\equiv 1 \pmod{3}$, is suggestive: the constant that
began the program as a threshold for digit matching may
also be connected to penetration into the critical strip.

% ============================================================
\section{Remarks}

\subsection*{The instrument}

The collision transform was introduced
in~\cite{paper16} to prove the Centered Collision PNT at
$s = 1$. The present paper observes that the same instrument
sees deeper. The digit function
$\delta(r) = \lfloor br/p \rfloor$, through its collision
statistics, produces a Dirichlet series whose convergence
appears, in the tested data, to extend into the critical
strip. The depth of this apparent region is a new
computational invariant of the collision structure.

\subsection*{Base as instrument}

The conventional view treats the positional base as a
notational choice with no mathematical content. This
program suggests otherwise. The base determines the bin
partition and the number of characters in the collision
transform, and at finite range the observed penetration
depth varies with the base. Different bases are different
instruments: base~$3$
has six characters and, in the present data, reaches
$s = 0.5$; base~$12$ has forty-eight and reaches only
$s = 0.7$ at the same computational depth. The base is
not mere notation here. It acts as the resolution of the
instrument that probes the prime.

\subsection*{The equivalence}

The digit function $\delta(r) = \lfloor br/p \rfloor$
encodes the prime's internal geometry. The collision
transform reads a finite stamp (the modular structure)
and isolates the aperiodic residual. The antisymmetry
restricts the residual to odd characters. The character
products have mixed signs, so convergence below $s = 1$
depends on cancellation between terms. Whether this
cancellation extends to $s = 1/2$, and how it relates
to the zeros of the relevant $L$-functions, remains open.

\subsection*{The arc}

The program began with the alignment of the prime~$3$
approaching $2/3$~\cite{paper1}. It built through the gate
width theorem~\cite{paper11}, the collision
transform~\cite{paper16}, and arrives here at the critical
strip.
Each step revealed structure or a new question: the gate
width is universal, the collision deviation is finitely
determined, the centered residual appears aperiodic in the
tested towers, and the problem of convergence at
$s = 1/2$ brings the program into direct contact with the
zero distribution of the relevant family of $L$-functions.
The inside of the prime and the zeros of the $L$-function
are connected through the collision transform, but the full
relationship remains open.

% ============================================================
\begin{thebibliography}{99}

\bibitem{davenport}
H.~Davenport,
\emph{Multiplicative Number Theory},
3rd~ed., Springer, 2000.

\bibitem{iwaniec-kowalski}
H.~Iwaniec and E.~Kowalski,
\emph{Analytic Number Theory},
AMS Colloquium Publications, vol.~53, 2004.

\bibitem{montgomery-vaughan}
H.~L.~Montgomery and R.~C.~Vaughan,
\emph{Multiplicative Number Theory I: Classical Theory},
Cambridge University Press, 2007.

\bibitem{paper1}
A.~S.~Petty,
Why the golden ratio selects the prime three,
research note, 2020.

\bibitem{paper11}
A.~S.~Petty,
Bin derangements and the gate width theorem,
research note, 2024.

\bibitem{paper15}
A.~S.~Petty,
The centered Collision PNT,
research note, 2025.

\bibitem{paper16}
A.~S.~Petty,
The collision transform,
research note, 2025.

\bibitem{nfield}
A.~S.~Petty,
\emph{nfield: A structural analysis engine for fractional field
invariants},
software.
\url{https://github.com/alexspetty/nfield}

\end{thebibliography}

\end{document}
